\documentclass[11pt,a4paper]{article}
\usepackage[T1]{fontenc}
\usepackage[utf8]{inputenc}
\usepackage{lmodern}
\usepackage[a4paper,left=27mm,right=27mm,top=26mm,bottom=28mm]{geometry}
\usepackage{microtype}
\usepackage{amsmath,amssymb,stmaryrd}
\usepackage{booktabs,tabularx,array}
\usepackage{xcolor}
\usepackage{titlesec}
\usepackage{fancyhdr}
\usepackage[hidelinks]{hyperref}

\definecolor{accent}{HTML}{1F4E79}
\definecolor{rulegray}{HTML}{9AA5B1}

\titleformat{\section}{\Large\bfseries\sffamily\color{accent}}{\thesection}{0.7em}{}
\titleformat{\subsection}{\large\bfseries\sffamily}{\thesubsection}{0.6em}{}
\titleformat{\paragraph}[runin]{\bfseries\sffamily}{}{0pt}{}[.]
\titlespacing*{\section}{0pt}{2.2ex plus .6ex}{1.0ex}
\titlespacing*{\subsection}{0pt}{1.8ex plus .4ex}{0.6ex}
\titlespacing*{\paragraph}{0pt}{1.2ex plus .3ex}{0.6em}

\setlength{\parindent}{0pt}
\setlength{\parskip}{0.62em plus 0.1em}
\renewcommand{\arraystretch}{1.25}
\newcolumntype{L}{>{\raggedright\arraybackslash}X}
\newcolumntype{P}[1]{>{\raggedright\arraybackslash}p{#1}}

\pagestyle{fancy}
\fancyhf{}
\renewcommand{\headrulewidth}{0pt}
\fancyfoot[L]{\footnotesize\sffamily\color{rulegray} ACT4E textbook -- progress 14 Oct -- 25 Nov 2026}
\fancyfoot[R]{\footnotesize\sffamily\color{rulegray}\thepage}

% notation
\newcommand{\cat}[1]{\mathbf{#1}}
\newcommand{\tL}[1]{\llparenthesis\mathbf{#1}\rrparenthesis}   % tuple version (product-like)
\newcommand{\tR}[1]{\rrparenthesis\mathbf{#1}\llparenthesis}   % indexed version (sum-like)
\newcommand{\Fb}{\mathrm{Fb}}
\newcommand{\Tr}{\mathrm{Tr}}
\newcommand{\St}{\mathbf{St}}
\newcommand{\R}{\mathbb{R}}
\newcommand{\file}[1]{{\small\texttt{#1}}}

\begin{document}

{\sffamily
{\color{accent}\rule{\textwidth}{1.2pt}}\\[0.9em]
{\LARGE\bfseries Textbook progress report}\\[0.35em]
{\large Applied Compositional Thinking for Engineers \quad$\cdot$\quad 14 October -- 25 November 2026}\\[0.5em]
{\small\color{rulegray} Prepared for the co-authors' meeting. Based on the repository history of
\texttt{JL-working-branch} from commit \texttt{764e99e} (9 Oct) to \texttt{463ab7d} (25 Nov).}\\[0.3em]
{\color{accent}\rule{\textwidth}{0.4pt}}
}

\section{Overview}

Over the last six weeks the work has moved through the book from front to back, in step with the
lecture schedule: the Actions chapter, then the first three chapters of the Categories part, the
Functors chapter, and finally both chapters of the Interconnections part. In the book sources
this amounts to about 60 commits (six merged pull requests, \#1111 and \#1113--\#1117), touching
111 files, with roughly 11{,}300 lines added and 9{,}400 removed. A good share of both numbers is
material that was relocated, so the three sections below describe what the numbers mean.

The three most important results, one per dimension, are these.

\paragraph{Content} The Interconnections part has gone from a skeleton to a written text. The
section on strict monoidal categories now carries a full stock of examples, there is a new section
on string diagrams, and the Feedback chapter has a new central section, ``Feedback and trace'',
with worked examples from relations, matrices, co-design, Mealy machines and LTI systems.

\paragraph{Structure} Dynamical systems and machines now follow one path through the book
(defined in Actions, composed in Modeling with categories, acting in a new Functors application
chapter, stacked in Parallel composition, looped in Feedback). The old chapter ``Application:
Dynamical system behaviors'' has been dissolved into this path. The tuple categories
$\tL{Set}$, $\tL{Rel}$, $\tL{Pos}$, $\tL{DP}$ have moved out of the Strictification section in
Higher structures and into Parallel composition, so that the Interconnections part can be read
without natural transformations.

\paragraph{Conceptual} Three decisions shape the new text: machines and LTI systems are defined
\emph{without} initial states; feedback operators and traces are treated together and told apart
by the \emph{delay} of an object; and every monoidal structure used for interconnection is given
in a strict model, of which there are now two families (product-like and sum-like).

\begin{table}[h]
\centering\small
\begin{tabularx}{\textwidth}{@{}P{26mm}P{47mm}L@{}}
\toprule
\textbf{When} & \textbf{Where} & \textbf{What}\\
\midrule
14--18 Oct & Algebra: Actions & Chapter re-sectioned; dynamical systems, machines, homomorphisms of actions written out\\
24--28 Oct & Categories: Intro on series composition; Categories & Intro chapter completed; definition of category rewritten; first examples\\
3--6 Nov & Categories: New categories from old & Subcategories, products, arrow, twisted arrow, slice and coslice reworked with applied examples\\
10--16 Nov & Functors; Modeling with categories & ``Functors generalize familiar concepts''; categories of LTI systems and of Moore and Mealy machines\\
18--19 Nov & Interconnections: Parallel composition & Strict monoidal examples, string diagrams, symmetric structure; systems chapter dissolved\\
23--25 Nov & Interconnections: Feedback; Functors applications & Initial states removed; ``Actions of dynamical systems and machines''; sum-like strict models; ``Feedback and trace''; duality section made strict\\
\bottomrule
\end{tabularx}
\caption{Where the work happened, in chronological order.}
\end{table}

\section{Content: what was written}

\subsection{Algebra part: Actions}

The chapter had five sections and now has eight; the former catch-all section ``Actions'' was
split up. The motivating examples (a robot on a grid, a
discrete-time system) are no longer development-only. Monoid, semigroup and group actions each have
their own section, with the definitions aligned to one another and to the uncurried form
$\mathrm{act}(m \fatsemi n, x) = \mathrm{act}(n, \mathrm{act}(m, x))$.

The section on dynamical systems grew from about 320 to 525 lines. New in it are a subsection on
sampling a continuous-time system (restriction of an action along the monoid morphism
$k \mapsto kh$, the sampled exponential, a remark on aliasing) and a rewritten subsection on LTI
systems: a discrete-time definition with a savings-account example, the continuous-time
definition, a mass between two springs, and proper systems.

``Machines'' is a new section of 515 lines: actions of sets (semiautomata, the transition monoid),
Moore machines, Mealy machines (a change detector; discrete-time LTI systems as Mealy machines), and
``More machines'', whose read-out returns a list of outputs (duplicator, discarder, terminator).

``Homomorphisms of actions'' grew from 42 to 423 lines. It now has examples (change of coordinates,
time-invariant filters, the 24-hour and the 12-hour clock, changing currency, conserved quantities)
and subsections on homomorphisms of discrete-time systems, continuous-time systems, Moore machines
and Mealy machines.

\subsection{Categories part}

\paragraph{Intro: Series composition} The sections ``Associativity'' and ``Neutral components''
were eight-line stubs and are now written (182 and 146 lines). ``Interfaces'' was extended so that
all three sections run the same cast of examples: electrical cables, sets and functions, preorders
and monotone functions, monoids, real matrices, co-design components, journeys, and currency
exchangers. ``Neutral components'' has subsections on examples with and without neutral
components, on adding them by convention (the cable of length zero), and on models where their
absence matters (exchangers with strictly positive commission).

\paragraph{Categories} The chapter text was rewritten. The section ``Categories'' builds up the
definition in four short subsections (objects and morphisms; composition and hom-sets;
associativity; identity morphisms), states it, and gives first examples: sets, real matrices, and
feasibility relations. Semicategories, isomorphisms, graphs and diagrams follow as separate
sections.

\paragraph{Modeling with categories} The category $\cat{LTI}$ is now defined here with real
coordinate spaces $\R^n$ as objects, with remarks on where the composition rule comes from and on
the state space $\R^{k+l}$ of a composite, a worked composition of two systems, and the definition
of equivalent LTI systems. A new section ``Machines'' shows that naive series composition of Moore
machines is associative only up to re-bracketing of states, defines the semicategory
$\cat{Moore}/{\sim}$ with a proof that it is one, explains why it has no identities, and then
defines the category $\cat{Mealy}/{\sim}$.

\paragraph{Categories of structures} New section ``Analysis'': Euclidean spaces with continuous,
smooth, and Lipschitz functions, each obtained by the same recipe (the property holds for
identities and is preserved by composition). The Algebra part gained the lemma that monoid
morphisms compose, which the intro chapter on series composition cites.

\paragraph{New categories from old} ``Subcategories'' was rewritten (subcategories of $\cat{Set}$
and $\cat{Rel}$, non-examples and near-misses, the relation to submonoids and free categories).
Products gained examples. The arrow construction has two extended examples, casting out nines and
fixed-point arithmetic. The twisted arrow construction now has a source-twisted and a
target-twisted version with a remark relating them. Slice and coslice have applied examples:
typed sets, typed graphs, all routes into an assembly, open circuits with terminals, and the
computation tree of a system.

\subsection{Functors part}

In the section ``Functors'' the standard constructions are now presented as an example followed
by a short graded exercise (``check that the proposed functor is indeed a functor''): cartesian
product, multiplication and exponentiation with a set, probability measures, differentiation,
fixed points, and LTI systems as Mealy machines.

The new section ``Functors generalize some familiar concepts'' (422 lines) covers graph
homomorphisms, monotone functions, monoid morphisms and monoid actions, each by unpacking the
definition of a functor and ending in a lemma. It continues with $\cat{Set}$-valued functors as
generalized actions, semicategory actions, contravariant actions, and the canonical actions of a
category (the hom-functors).

The new section ``Actions of dynamical systems and machines'' (400 lines), in the chapter
``Applications: some extended examples'', defines the categories $\cat{LTI}_*$ and
$\cat{Mealy}_*/{\sim}$ of systems with a chosen initial state and lets them act: LTI systems on
continuous-time signals and on discrete-time streams, Mealy and Moore machines on streams. It
includes the lemma that equivalent LTI systems act in the same way and a remark that the converse
fails.

\subsection{Interconnections part}

\paragraph{Parallel composition} The section ``Strict monoidal categories'' went from 150 to
about 2{,}360 lines. Roughly half of that is the tuple-category material brought over from
Strictification and reworked; the rest is new. The examples are now organized as follows: real
matrices with the dimension-adding product and with the Kronecker product; categories whose
objects are natural numbers (functions between finite sets, polynomial maps, signal-flow
diagrams); products ``like the cartesian product'' ($\tL{Set}$, $\tL{Pre}$, $\tL{Pos}$,
$\tL{Rel}$, $\tL{mRel}$, $\tL{Bool\text{-}Prof}$); products ``like the sum'' ($\tR{Set}$,
$\tR{Rel}$, $\tR{Pre}$, $\tR{Pos}$, built from indexed elements); applied examples (co-design,
Mealy machines, LTI systems); monoidal preorders; and one-object strict monoidal categories, with
the Eckmann--Hilton argument that these are commutative monoids.

\paragraph{String diagrams} The section ``The graphical calculus of string diagrams'' is new (735
lines) and part of the public text: boxes and wires, why the pictures are unambiguous, sliding
boxes, from pictures to formulas, calculating with pictures (the Joyal--Street theorem, stated
informally), and a worked example on adding matrices.

\paragraph{Symmetric strict monoidal categories} This section went from 47 to 348 lines. It now
gives the braiding in each running example: LTI systems up to equivalence, tuples, relations
between tuples, indexed elements, Mealy machines, and matrices with the Kronecker product.

\paragraph{Feedback} ``Feedback and trace'' (1{,}314 lines) replaces two short sections of 174 and
98 lines. It defines feedback operators, traces, and the delay of a feedback operator, with a
figure for each axiom. The basic examples are relations with the product-like and with the
sum-like monoidal product, monotone relations, Boolean profunctors, matrices with the Kronecker
product, matrices with the direct sum (as a non-example), and linear relations. The applied
examples are co-design (a vehicle that transports its own battery; battery modules) and the
construction of stateful morphisms, with Mealy machines and LTI systems as its two instances.

\paragraph{Canonical trace, duality} The section was rewritten for the strict setting: snake
equations without associators and unitors, dual morphisms, and the canonical trace. New examples:
$\tL{mRel}$ is compact closed, three categories that are not, the dual of a matrix is its
transpose, and the canonical trace for monotone relations. The four graded exercises of the
section were restated for the strict categories.

\paragraph{Supporting material elsewhere} The chapter on monotone relations gained Boolean
profunctors and composition as Boolean matrix multiplication, which the interconnection examples
use.

\section{Structure: how the book was reorganized}

\begin{table}[h]
\centering\small
\begin{tabularx}{\textwidth}{@{}P{33mm}LL@{}}
\toprule
\textbf{Place} & \textbf{9 October} & \textbf{25 November}\\
\midrule
Actions chapter & ``Actions'' (monoid, semigroup and set actions in one section), Vector spaces, Dynamical systems (including Moore and Mealy machines), Homomorphisms, Group actions
 & Motivations, Monoid actions, Semigroup actions, Group actions, Modules and vector spaces, Dynamical systems, Machines, Homomorphisms of actions\\
Categories chapter & Semicategory defined first; category as a semicategory with identities
 & Category defined first; semicategories, isomorphisms, graphs and diagrams as separate sections\\
Functors part & After the intro chapter: ``Functors'', ``Functors at work'' (development-only content), ``Application: Dynamical system behaviors''
 & After the intro chapter: ``Functors'' and ``Applications: some extended examples''\\
Parallel composition & Graphical calculus development-only; ``Props'' a section
 & Graphical calculus public; ``Props'' a subsection of the symmetric section\\
Tuple categories & In Higher structures, under Strictification, as ``practical strictifications''
 & In Parallel composition, as the main examples of strict monoidal categories\\
Feedback chapter & ``Feedback categories'' and ``Traced symmetric strict monoidal categories''
 & One section ``Feedback and trace''\\
\bottomrule
\end{tabularx}
\caption{The main structural changes.}
\end{table}

\paragraph{One path for systems} Before, LTI systems and machines were introduced in the Actions
chapter and then treated again, largely from scratch, in the chapter ``Application: Dynamical
system behaviors'' at the end of the Functors part. That chapter has been taken apart. Its
definitions went to where the concepts first appear: More machines and the modelling remarks on
LTI systems to Actions; the category of LTI systems, the exercise that it is a category, and
equivalence of systems to Modeling with categories; semicategory actions to the Functors chapter;
and the actions on signals and streams to the new application section. A system now appears five
times, each time with one more layer of structure: as a single object (Actions), as a morphism
(Modeling with categories), as something that acts (Functors applications), as something that can
be stacked (Parallel composition), and as something whose loop can be closed (Feedback).

\paragraph{Strict first} The plan discussed in August, a strict-only Interconnections part with
the non-strict theory deferred to Higher structures, is now largely carried out in the text. The
tuple categories were moved out of ``Strictification'' and are introduced in Parallel composition
on their own terms. The relation examples for traces were cut from the non-strict chapter on
traced monoidal categories and rewritten strictly in ``Feedback and trace''. The duality section
no longer uses associators or unitors. The $\tL{Set}$ version of Moore machines moved to the
appendix on parallel composition for semicategories. What remains in Higher structures under
``Strictification'' is the motivation and the statement of Mac Lane's theorem.

\paragraph{Within-chapter order} In Categories of structures, linear algebra now comes before
actions. In the Functors chapter, the three ``functors generalize'' stubs that sat at the end of
the first section became a section of their own, placed before the further examples.

\paragraph{Patterns applied throughout} All new definitions use the Constituents and Conditions
layout, with identities and unitality listed last. Results of the form ``$X$ is a functor'' follow
the example-plus-graded-exercise pattern, which keeps the exercise names and labels that the
homework sheets and solution files refer to. The instructors' index of graded exercises was
cleaned up (stale entries and point counts removed), and broken references were fixed after each
move.

\section{Conceptual: what changed in how the material is presented}

\paragraph{Series composition before categories} The intro chapter now isolates three patterns
(interfaces, associativity, neutral components) across one cast of examples, and the definition
of a category is presented as the name for these patterns. Neutral components are treated as a
modelling choice: some models have them, some gain them by convention, and in some their absence
carries meaning. This gives semicategories a motivation of their own, as the notion for models
without neutral components, while the category is the primary definition. The analogy
semigroup is to monoid as semicategory is to category is used consistently, down to the order of
items in the definitions.

\paragraph{Dynamical systems versus machines} The Actions chapter now separates systems driven by
time (the acting monoid is $\mathbb{N}$ or $\R$, its elements are durations) from systems driven
by events (the acting monoid is a monoid of lists, only the order matters). The two overlap in
discrete time, and the text says so in a remark titled ``Same structure, two readings'': a
discrete-time system is a machine with one kind of event, and a discrete-time LTI system is a
Mealy machine.

\paragraph{No initial states} Moore, Mealy and More machines and LTI systems are now defined
without an initial state; the initial state ``is not part of the machine: it is a choice that we make when we run the
machine''. The
reason comes from feedback: closing a loop turns the fed-back wire into part of the state, and a
definition with initial states would then have to invent an initial value for that wire. The
consequence is a clean split between specification and behaviour. The categories $\cat{LTI}$,
$\cat{Moore}/{\sim}$ and $\cat{Mealy}/{\sim}$ record how systems are specified and connected.
Behaviour belongs to a system that has been started, so the actions on signals and streams are
functors from the categories $\cat{LTI}_*$ and $\cat{Mealy}_*/{\sim}$ of systems with an initial
state. A remark explains why the choice of initial state cannot be left to the functor.

\paragraph{Behaviour is an action} The Functors chapter now states that a functor to $\cat{Set}$
is an action of a category, in the same sense in which a monoid action is a functor from a
one-object category. With this in place, ``what a system does'' has a precise form, and
functoriality says that the behaviour of a series composition is the composition of the
behaviours. In the course of this work a lemma in the old systems chapter turned out to be false:
the ``non-standard'' action of Moore machines is not compatible with series composition (two
delays composed give a delay by two steps, while the action of the composite gives a delay by
one). The lemma and its exercise were removed from the book.

\paragraph{Equality up to isomorphism, made explicit} Several constructions hold only after
passing to a quotient, and the text now says where and why. Moore and Mealy machines compose
associatively only up to re-bracketing of states. Stacking of LTI systems is not functorial on
the nose, because the two sides of the interchange law list the state coordinates in different
orders; it is functorial for systems up to equivalence. The same phenomenon reappears for the
braiding of LTI systems and for stateful morphisms.

\paragraph{Two families of strict models} The cartesian product and the sum of sets are associative and unital only up to renaming, so
neither gives a strict monoidal category directly, and the text repairs each with a concrete
model. For product-like structures the model is flat tuples: an element of a product is an
element of the first object \emph{and} one of the second, and something travels along every
wire. For sum-like structures the model is new: sets of indexed elements, where an element of a
product is an element of the first object \emph{or} of the second, together with an index, and
the wires are alternatives. Matrices provide the contrasting case of one category with two
strict monoidal products (adding dimensions, and Kronecker).

\paragraph{Feedback versus trace} A trace is now defined as a feedback operator that also
satisfies dinaturality for all morphisms and yanking. The new notion of the delay of an object,
$\partial_Z := \Fb^Z_{Z,Z}(\sigma_{Z,Z})$, describes what happens to something that travels once
around a loop. Yanking says that all delays are identities, so one non-identity delay shows that
a feedback operator is not a trace. The text follows a naming rule: operators that are traces are
written $\Tr$, the others $\Fb$, and for each of the latter the text says why it is not a trace.

\begin{table}[h]
\centering\small
\begin{tabularx}{\textwidth}{@{}P{43mm}LP{20mm}P{33mm}@{}}
\toprule
\textbf{Category} & \textbf{Closing a loop means} & \textbf{Status} & \textbf{Delay}\\
\midrule
$\tL{Rel}$, $\tL{mRel}$, $\tL{Bool\text{-}Prof}$, $\tL{Feas}$, $\tL{DP}$ & there is a loop value $c$ that is consistent on both sides & trace & identity\\
$\tR{Rel}$ & iterate: go around the loop any number of times & trace & identity\\
Matrices, Kronecker product & block-wise linear-algebra trace & trace & identity\\
Linear relations & same formula as for $\tL{Rel}$ & trace & identity\\
Matrices, direct sum & $M_{11} + M_{12}(1 - M_{22})^{-1}M_{21}$, not always defined & neither & --\\
$\St(\cat{C})$, stateful morphisms & the wire becomes part of the state & feedback only & stores its input for one step\\
$\tL{Mealy/{\sim}} = \St(\tL{Set})$ & as above & feedback only & one-step register\\
$\cat{LTI}/{\sim} = \St(\cat{Mat}_\R)$ & as above & feedback only & integrator\\
\bottomrule
\end{tabularx}
\caption{The feedback examples as they now stand in ``Feedback and trace''.}
\end{table}

\paragraph{Stateful morphisms unify machines and LTI systems} The construction $\St(\cat{C})$
(after Katis, Sabadini and Walters) equips the morphisms of a symmetric strict monoidal category
with a state object. Applied to $\tL{Set}$ it gives Mealy machines up to isomorphism; applied to
real matrices it gives LTI systems up to equivalence, an LTI system being the block matrix
$\bigl[\begin{smallmatrix} D & C\\ B & A\end{smallmatrix}\bigr]$. The feedback operator for LTI
systems is therefore now in the book, with an explicit formula. The text points out what it
means: the loop is closed through an integrator ($u_2 = w$, $\dot w = y_2$), which is different
from the direct connection $u_2 = y_2$ that the partial trace in the appendix describes.

\paragraph{What duality adds} In the strict duality section, the canonical trace of $\tL{mRel}$
comes out in a ``$\leq$'' form, so the co-design reading of a loop (the requirement fed back must
not exceed the functionality provided) is built into evaluation and coevaluation. The example of
categories that are not compact closed shows two further points: for matrices with the direct sum
only $\R^0$ is dualizable, which matches the missing trace there, and $\tR{Rel}$ has a trace
although only the empty set is dualizable, so a trace need not come from duals.

\paragraph{Groundwork not yet in the book} Four standalone notes with full proofs were written
alongside the text and sit in the folder \file{Claude outputs}: on Mealy feedback with separate
dynamics and read-out; on the LTI feedback operator and its relation to the partial trace; on
feedback operators for matrices; and a longer note showing that co-design queries form the
Kleisli category of the upper-set monad, isomorphic to monotone relations compatibly with series
composition, parallel composition and feedback. The last one contains a proposed map for
integrating the result into the book and a list of errata in the draft chapters on querying.

\vspace{1.2em}
{\color{rulegray}\rule{\textwidth}{0.4pt}}

{\footnotesize\color{black!65}
\textbf{Scope and method.} This report covers the textbook sources under \file{volumes/vol1}
together with the shared macros and bibliography. The slides book, homework sheets, quizzes and
exams were updated in parallel during the same period and are not described here. Line counts are
raw counts of the \LaTeX{} source files, including comments and figure code, and are meant only to
indicate scale. Section and chapter names are quoted as they appear in the sources on
25 November 2026.}

\end{document}
